% Pursuing Stacks (cotes 134-2 to 134-8): archive cote 134-4, site page 10 (PDF page 11), typescript page 98.
% typescript. Draft compared: ps3.tex:479-551.
% AI-transcribed by Claude Opus 5.5 (claude-opus-5-5, adaptive thinking, effort high), Message Batches API, batch ps1 (msgbatch_01Rv5hCaNTtUt4JHx1yqoHPL), 2026-11-12; not yet checked by eye.
\scanpage{134-4/10}
\mspage{(98)}
tractible'' as synonymous to ``aspheric'', as \xout{\unsure{for the}} in the context of
topological spaces, or CW spaces at any rate, the two notions do indeed
coincide, finally I came to use \inserted{rather} the word ``aspheric'' systematically, as it
fits nicely with the notion of an aspheric morphism of topoi...). THe
most evident example \deleted{of all} of contractible categories \inserted[typed]{are} \xout{xx} the categories
with final object. Thus I have the strong feeling that the condition of
contractibility of the \xout{\illegible{1}} objects $i(a)$ in (Cat) is ``the right'' genera-%
lization of the \xout{condition}assumption made in the ``key result'', namely
that the $i(a)$\hand{'}s have final element\hand{s}. Also, it seems now likely that the
numerous cases of statements, when to check some property for arbitrary
$C$, it turned out to be enough to check it for \xout{contractible} $C$ with a final
object, may well generalize to more general cases, with (Cat) replaced
by some $M$ and the reduction is from arbitrary $C$ in $M$ to contractible ones.

The next thing to do is to develop a little the notion of contractibi%
lity of objects and of homotopies between maps, and to get the criterion
just announced for (3) to be a weak equivalence for any $C$. After this,
handling the question of the first map in (1)
\[ \text{(4)} \qquad A_{/C} \hand{\to} A_{//C} \]
being a weak equivalence for any $C$, in terms of the asphericity criteri-%
on for a functor, will turn out pretty much formal, and we'll finally
be able to state a new version of the ``key result'' about test functors
$i$: $A \hand{\to}$ (\unsure{Cat}), with \inserted{this time} twice as many equivalent formulations of the same
property. On n'arrête pas le progrès !

\hand{\unsure{Mer.}}
31.3.

48. It's time now to develop some generalities about homotopy classes
of maps, the relation of homotopy between objects of a modelizer, and th%
corresponding notion of contractibility. For the time being, it will be
enough to start with any category $M$, endowed with a set $W \hand{\subset} \mathrm{Fl}(M)$ of
arrows (the``weak equivalences''), satisfying the mild saturation assumpti-%
ons a)b)c) of par.37 (p.59). On $M$ we'll assume for the time being that
there exists (at least one) homotopy interval\deleted{\unsure{,}}\inserted{$\underline{I}=$}$(I,\hand{\delta}_0,\hand{\delta}_1)$ in $M$ (loc. cit.)
which implies also that $M$ has a final object, which I denote by $e_M$ or
simply \hand{$e$} , and equally an initial element $\varnothing_M$ . I'm not too sure yet whe-%
ther we'll really need that the latter be a strict initial element, as
required in the definition of a homotopy interval on page 59 (it was used
in the generalities of pages 59 and 60 only for the corollary on the Law-%
vere element ...). Whether or not will appear soon enough ! I'll not as-%
sume it till I am forced too. \hand{To be safe, we'll assume on the other
hand that $M$ admits binary products.}
\margin{1°}
